\section{Submission Package Checklist and Outstanding TODOs}
\label{apx:submission-checklist}

This appendix records the items that are still missing or still need to be reformatted for a complete NGCC submission package.

\subsection{Algorithm Text TODOs}

\begin{itemize}
    \item A submission-style security-claim table for the selected NGCC 128/256/512-bit families has now been added in \cref{tab:ngcc-parameters}. \todo{If desired, add a supplementary note explaining how the single-key $q=2^{64}$ tuning used for parameter selection relates to the evaluation document's broader ``at most $2^{80}$ queries'' convention.}
    \item The recommended families now use $h=70$ and are tuned with $q=2^{64}$, which supports the call's requirement of at least $2^{64}$ signatures under one key pair in the intended single-key model. \todo{State explicitly the supported message-length bound and mirror the final signing-capacity wording in the finalized submission forms.}
    \item \todo{Add a short, explicit correctness statement in the exact notation and wording preferred for the submission package.}
    \item \todo{Expand the discussion of side-channel resistance, multi-key security, and repeated-key usage with either concrete evidence or a clearly scoped limitation statement.}
    \item \todo{Add optional benchmark data for 32-bit embedded platforms and hardware platforms if these results become available.}
    \item \todo{Add quantitative memory-consumption measurements instead of the current qualitative remark on low stack usage.}
\end{itemize}

\subsection{Required Forms and Statements}

\begin{itemize}
    \item \todo{Provide the signed algorithm basic-information form for all submitters and the designated contact person.}
    \item \todo{Provide the signed submitter intellectual-property statement.}
    \item \todo{Provide signed patent-owner / patent-applicant statements for all relevant rights, if any.}
    \item \todo{Provide the signed ownership statement for the reference and optimized implementations.}
\end{itemize}

\subsection{Implementation Package}

\begin{itemize}
    \item First-stage ISO C reference implementations are now present for \texttt{Kylin-128}, \texttt{Kylin-256}, and \texttt{Kylin-512}.
    \item First-stage optimized 64-bit PC implementations are now present for all three instances; they currently share the portable code path with an optimized build profile. \todo{Replace or supplement this with deeper architecture-specific optimization before final submission, if desired.}
    \item The implementations use the NGCC SIG API provided in \texttt{API\_PKC.zip}.
    \item A repository-local \texttt{README} for the implementation package has been added.
    \item \todo{Add final benchmark tables and memory-consumption measurements for the completed implementation package.}
\end{itemize}

\subsection{Test Vectors}

\begin{itemize}
    \item Known-answer test vectors for all three submitted instances have been generated using the NGCC \texttt{KAT\_SIG.c} driver.
    \item The generated files are \texttt{KAT\_SIG\_Kylin-128.txt}, \texttt{KAT\_SIG\_Kylin-256.txt}, and \texttt{KAT\_SIG\_Kylin-512.txt}.
    \item The reference build scripts regenerate and copy the published test vectors with \texttt{make kat}.
\end{itemize}
